\section{Technical Appendix \& Deep Dive}

\begin{frame}[label=backup_rootsystem]{Definition 4.4.1: Root System}
    Let $E$ be a finite-dimensional Euclidean vector space, with $(\cdot, \cdot) : E \times E \to \mathbb{R}$ the scalar product. A root system $\mathcal{R} \subset E$ is a finite set which satisfies:
    \begin{itemize}
        \item[a)] The $\mathbb{R}$-span of $\mathcal{R}$ is $E$.
        \item[b)] For $\alpha \in \mathcal{R}$, if $\lambda \alpha \in \mathcal{R}$ for $\lambda \in \mathbb{R}$, then $\lambda = \pm 1$.
        \item[c)] \textbf{Crystallographic condition:} $\forall \alpha, \beta \in \mathcal{R}: \quad n_{\alpha\beta} := \frac{2(\alpha, \beta)}{(\alpha, \alpha)} \in \mathbb{Z}$.
        \item[d)] \textbf{Reflection invariance:} $\forall \alpha, \beta \in \mathcal{R}: \quad s_\alpha(\beta) := \beta - n_{\alpha\beta}\alpha \in \mathcal{R}$.
    \end{itemize}
\end{frame}

\begin{frame}[label=backup_angles]{Extended Proof: Angle Quantization}
    From Definition 4.4.1 (c), we know $n_{\alpha\beta} \in \mathbb{Z}$.
    Using Euclidean geometry definitions:
    \begin{align*}
        \|\alpha\|                & = \sqrt{(\alpha, \alpha)}                      \\
        \cos \theta_{\alpha\beta} & = \frac{(\alpha, \beta)}{\|\alpha\| \|\beta\|}
    \end{align*}

    Substitute $(\alpha, \beta)$ into the product of integers:
    \begin{align*}
        n_{\alpha\beta} n_{\beta\alpha} & = \left( 2 \cos \theta_{\alpha\beta} \frac{\|\beta\|}{\|\alpha\|} \right) \left( 2 \cos \theta_{\alpha\beta} \frac{\|\alpha\|}{\|\beta\|} \right) \\
                                        & = 4 \cos^2 \theta_{\alpha\beta}
    \end{align*}
    Since $\cos^2 \theta \in [0,1]$, the integer product must be $\in \{0,1,2,3,4\}$.
\end{frame}

\begin{frame}[label=backup_astring]{Theorem 4.4.2: Root Strings}
    For $\alpha, \beta \ne \alpha$ two roots, the $\alpha$-string through $\beta$ (the set $\beta - p\alpha, \dots, \beta + q\alpha$) satisfies:
    \begin{itemize}
        \item[a)] The $\alpha$-string through $\alpha$ is $\{-\alpha, 0, \alpha\}$.
        \item[b)] $p - q = n_{\alpha\beta} \equiv \frac{2(\alpha, \beta)}{(\alpha, \alpha)}$.
        \item[c)] For $\alpha \ne \pm \beta$, the string has at most 4 elements ($p+q \le 3$).
        \item[d)] \textbf{Crucial property for drawings:} If $(\alpha, \beta) > 0$, then $\alpha - \beta \in \mathcal{R}$. If $(\alpha, \beta) < 0$, then $\alpha + \beta \in \mathcal{R}$.
    \end{itemize}
\end{frame}

\begin{frame}[label=backup_simpleroots]{Definition 4.4.3: Positive \& Simple Roots}
    Let $t \in E$ be such that $(t, \alpha) \ne 0$ for all $\alpha \in \mathcal{R}$ (always possible because $\mathcal{R}$ is finite).

    \begin{itemize}
        \item A root $\alpha \in \mathcal{R}$ is called \textbf{positive} ($\alpha \in \mathcal{R}^+$) if $(t, \alpha) > 0$, and \textbf{negative} if $(t, \alpha) < 0$.
        \item A positive root is called \textbf{simple} if it cannot be written as the sum of two positive roots.
    \end{itemize}

    The set of simple roots is denoted $\Delta \subset \mathcal{R}$.

    \vspace{0.2cm}
    \textbf{Lemma 4.4.4:} If $\alpha, \beta \in \Delta$, then:
    a) $\alpha - \beta$ and $\beta - \alpha$ are not roots.
    b) $(\alpha, \beta) \le 0$ (obtuse angles only).
    c) $\Delta$ is linearly independent.
\end{frame}